---
title: "श्री हनुमान चालीसा"
geometry: "paperwidth=8.5cm, paperheight=5.4cm, margin=0.3cm"
fontsize: 6pt
documentclass: article
header-includes: |
  \usepackage[utf8]{inputenc}
  \usepackage{xcolor}
  \usepackage{fancyhdr}
  \pagestyle{fancy}
  \fancyhf{}
  \renewcommand{\headrulewidth}{0pt}
  \setlength{\parindent}{0pt}
  \setlength{\parskip}{1pt}
---

\begin{center}
{\Large \textbf{श्री हनुमान चालीसा}}
\end{center}

\vspace{2mm}

\textbf{दोहा}

श्रीगुरु चरन सरोज रज निज मनु मुकुरु सुधारि ।\\
बरनउँ रघुबर बिमल जसु जो दायकु फल चारि ॥

बुद्धिहीन तनु जानिके, सुमिरौं पवन कुमार ।\\
बल बुधि विद्या देहु मोहि, हरहु कलेश विकार ॥

\vspace{1mm}

\textbf{चौपाई}

जय हनुमान ज्ञान गुन सागर । जय कपीस तिहुँ लोक उजागर ॥१॥\\
राम दूत अतुलित बल धामा । अंजनि पुत्र पवनसुत नामा ॥२॥\\
महाबीर बिक्रम बजरंगी । कुमति निवार सुमति के संगी ॥३॥\\
कंचन बरन बिराज सुबेसा । कानन कुंडल कुंचित केसा ॥४॥\\
हाथ बज्र औ ध्वजा बिराजे । काँधे मूँज जनेऊ साजे ॥५॥

शंकर सुवन केसरी नंदन । तेज प्रताप महा जग बंदन ॥६॥\\
विद्यावान गुनी अति चातुर । राम काज करिबे को आतुर ॥७॥\\
प्रभु चरित्र सुनिबे को रसिया । राम लखन सीता मन बसिया ॥८॥\\
सूक्ष्म रूप धरि सियहिं दिखावा । विकट रूप धरि लंक जरावा ॥९॥\\
भीम रूप धरि असुर संहारे । रामचंद्र के काज सवाँरे ॥१०॥

लाय सजीवन लखन जियाए । श्री रघुबीर हरषि उर लाए ॥११॥\\
रघुपति कीन्ही बहुत बड़ाई । तुम मम प्रिय भरत-हि सम भाई ॥१२॥\\
सहस बदन तुम्हरो जस गावै । अस कहि श्रीपति कंठ लगावै ॥१३॥\\
सनकादिक ब्रह्मादि मुनीसा । नारद सारद सहित अहीसा ॥१४॥\\
जम कुबेर दिगपाल जहाँ ते । कवि कोविद कहि सके कहाँ ते ॥१५॥

तुम उपकार सुग्रीवहिं कीन्हा । राम मिलाय राज पद दीन्हा ॥१६॥\\
तुम्हरो मंत्र विभीषण माना । लंकेश्वर भए सब जग जाना ॥१७॥\\
जुग सहस्त्र जोजन पर भानू । लील्यो ताहि मधुर फल जानू ॥१८॥\\
प्रभु मुद्रिका मेलि मुख माही । जलधि लाँघि गए अचरज नाही ॥१९॥\\
दुर्गम काज जगत के जेते । सुगम अनुग्रह तुम्हरे तेते ॥२०॥

राम दुआरे तुम रखवारे । होत न आज्ञा बिनु पैसारे ॥२१॥\\
सब सुख लहै तुम्हारी सरना । तुम रक्षक काहू को डरना ॥२२॥\\
आपन तेज समहारो आपै । तीनों लोक हाँक ते काँपै ॥२३॥\\
भूत पिसाच निकट नहिं आवै । महाबीर जब नाम सुनावै ॥२४॥\\
नासै रोग हरै सब पीरा । जपत निरंतर हनुमत बीरा ॥२५॥

संकट ते हनुमान छुड़ावै । मन क्रम बचन ध्यान जो लावै ॥२६॥\\
सब पर राम तपस्वी राजा । तिन के काज सकल तुम साजा ॥२७॥\\
और मनोरथ जो कोई लावै । सोई अमित जीवन फल पावै ॥२८॥\\
चारो जुग परताप तुम्हारा । है परसिद्ध जगत उजियारा ॥२९॥\\
साधु संत के तुम रखवारे । असुर निकंदन राम दुलारे ॥३०॥

अष्ट सिद्धि नव निधि के दाता । अस बर दीन्ह जानकी माता ॥३१॥\\
राम रसायन तुम्हरे पासा । सदा रहो रघुपति के दासा ॥३२॥\\
तुम्हरे भजन राम को पावै । जनम जनम के दुख बिसरावै ॥३३॥\\
अंत काल रघुबर पुर जाई । जहाँ जन्म हरि भक्त कहाई ॥३४॥\\
और देवता चित्त न धरई । हनुमत सेइ सर्व सुख करई ॥३५॥

संकट कटै मिटै सब पीरा । जो सुमिरै हनुमत बलबीरा ॥३६॥\\
जै जै जै हनुमान गोसाईं । कृपा करहु गुरुदेव की नाईं ॥३७॥\\
जो सत बार पाठ कर कोई । छूटहि बंदि महा सुख होई ॥३८॥\\
जो यह पढ़ै हनुमान चालीसा । होय सिद्धि साखी गौरीसा ॥३९॥\\
तुलसीदास सदा हरि चेरा । कीजै नाथ हृदय महँ डेरा ॥४०॥

\vspace{1mm}

\textbf{दोहा}

पवनतनय संकट हरन, मंगल मूरति रूप ।\\
राम लखन सीता सहित, हृदय बसहु सुर भूप ॥

\vspace{2mm}

\begin{center}
{\small \textbf{॥ सीया वर रामचंद्र की जय ॥}}\\
{\small \textbf{॥ पवन सुत हनुमान की जय ॥}}\\
{\small \textbf{॥ उमापति महादेव की जय ॥}}
\end{center}

\newpage

\begin{center}
{\Large \textbf{हनुमान आरती}}
\end{center}

\vspace{2mm}

आरती कीजै हनुमान लला की ।\\
दुष्ट दलन रघुनाथ कला की ॥

जाके बल से गिरिवर काँपे ।\\
रोग दोष जाके निकट न झाँके ॥

अंजनी गर्भ संभव सुर भूपा ।\\
कष्ट निवारण मंगल रूपा ॥

लाल देह लाली लसे ।\\
ताज देख राघव गुण गासे ॥

हाथ बज्र औ ध्वजा विराजै ।\\
काँधे मूंज जनेऊ साजे ॥

सुर नर मुनि आरती उतारें ।\\
जय जय जय हनुमान उचारें ॥

कंकन किंकिनी नूपुर धुनि सोहै ।\\
ग्वाल बाल नृत्य सुख होहै ॥

सुमिरत हि ते सकल दुख भागे ।\\
रोग शोक संताप न लागे ॥

\vspace{3mm}

\begin{center}
{\footnotesize Print this on sturdy paper, cut to credit card size (8.5cm × 5.4cm)}\\
{\footnotesize Laminate for durability and wallet storage}
\end{center}